\documentclass[12pt, a4paper]{article}
\usepackage[T1]{fontenc}
\usepackage{lmodern}
\usepackage{amsmath, amsthm, amssymb, mathtools}
\usepackage[margin=1in]{geometry}
\usepackage{xr-hyper}
\usepackage[colorlinks=true,linkcolor=blue,citecolor=blue]{hyperref}
\usepackage{cleveref}
\usepackage{graphicx, subcaption}
\usepackage{natbib}
\usepackage{booktabs}
\usepackage{enumerate}

% Code references: scaled monospace with line-breaking
\newcommand{\code}[1]{\texttt{\small #1}}

\newtheorem{theorem}{Theorem}[section]
\newtheorem{lemma}[theorem]{Lemma}
\newtheorem{proposition}[theorem]{Proposition}
\newtheorem{corollary}[theorem]{Corollary}
\newtheorem{definition}{Definition}[section]
\newtheorem{remark}{Remark}[section]
\newtheorem{conjecture}[theorem]{Conjecture}
\newtheorem{derivation}[theorem]{Derivation}
\newtheorem{observation}[theorem]{Observation}
\numberwithin{equation}{section}

\newcommand{\dd}{\mathrm{d}}
\newcommand{\seriestitle}{Why Rotating Fluids Make Polygons}


\externaldocument[I-]{../paper-1-mathematics/main}

\title{\seriestitle\\[6pt]\large Supplement A: Number-Theoretic Aspects of the Palindromic Hierarchy}
\author{Gordon Speagle}
\date{}

\begin{document}
\maketitle

\begin{abstract}
This supplement develops the number-theoretic consequences of the palindromic
stability hierarchy established in Paper~I.
The Bolza modular form $\eta(8z)\eta(16z)$ (level~128, weight~1)
is identified from the palindromic quartic; its Hecke eigenvalues
follow a Pythagorean sign rule.
The Fibonacci systole tower connects palindromic thresholds to
congruence surfaces.
The GPY sieve ratio~$1.92$ for ``Bolza twin primes'' provides
an analytic-number-theory fingerprint.
The vortex framework identifies concrete Langlands test cases:
non-solvable palindromic polynomials with $S_5$ Galois group
give four-dimensional Artin representations of $\mathrm{GL}(4)$
for which no automorphic lift is known.
\end{abstract}


%% ====================================================================
%% Extracted from Paper I, §6.5
%% ====================================================================

\section{Vortex dynamics as a Langlands laboratory}
\label{sec:langlands-lab}

The Frobenius data at a prime~$p$---the cycle type of
$\mathrm{Frob}_p$ in the Galois group of the palindromic
polynomial---is also the \emph{stability fingerprint} of~$p$:
the number of roots of~$Q(\xi)$ modulo~$p$ determines
which $\mathbb{Z}_N$ modes are visible on the surface
$\Gamma(p)\backslash\mathbf{H}^2$
(Proposition~\ref{prop:legendre-selection}).
Langlands reciprocity predicts this Frobenius data
matches an automorphic form on~$\mathrm{GL}(4)/\mathbb{Q}$.

\paragraph{Chebotarev test.}
For the three $S_5$ palindromic polynomials of
Proposition~\ref{prop:nonsolvable}, the root count
distribution over $428$ primes~$p < 3000$ matches the
Chebotarev prediction to within statistical error:
\begin{center}
\begin{tabular}{ccc}
\toprule
$\#\{\text{roots of } Q \bmod p\}$ &
  Chebotarev & Observed (mean) \\
\midrule
$0$ & $60.7\%$ & $60.9\%$ \\
$2$ & $30.3\%$ & $32.0\%$ \\
$4$ & $7.6\%$  & $5.9\%$ \\
$6$ & $1.3\%$  & $1.1\%$ \\
\bottomrule
\end{tabular}
\end{center}
The palindromic constraint (roots in pairs $\xi, 1/\xi$)
forces the root count to be even at every unramified prime,
confirmed with zero violations across $1284$ prime--polynomial
pairs.

\paragraph{What the vortex framework contributes.}
The framework provides three levels of evidence:
\begin{enumerate}
  \item \emph{Concrete test cases.}
    The LMFDB Artin representations
    \texttt{4.38569.5t5.a.a}, \texttt{4.65657.5t5.a.a},
    \texttt{4.24217.5t5.a.a}
    (Remark~\ref{rem:langlands-test}) are specific
    4-dimensional representations of $S_5$ for which no
    automorphic lift is known.
    The Frobenius data is now computable from the vortex
    stability fingerprint at each prime.
  \item \emph{A physical constraint.}
    The palindromic symmetry $Q(\xi) = \xi^{10}\,Q(1/\xi)$
    forces the functional equation of the Artin $L$-function.
    This is a \emph{necessary} condition for the existence of
    the automorphic lift---and it is guaranteed by the conformal
    inversion symmetry of the Green's function
    (Section~\ref{sec:dilation-invariance}), not assumed.
  \item \emph{The KK product at threshold.}
    Near a palindromic threshold, the spectral projection
    $P_- \in L(\Gamma)$ is non-trivial
    ($\tau(P_-) \approx 0.61$ on the Bolza surface).
    If the automorphic lift $\pi$ exists, it constrains
    $\tau(P_-)$ to lie in $\frac{1}{|\Gamma|}\mathbb{Z}$
    (by the rationality of $L^2$-Betti numbers).
    A high-precision computation of $\tau(P_-)$ would test
    this rationality prediction.
\end{enumerate}

\paragraph{What it does not provide.}
The palindromic symmetry gives the functional equation, and
the Chebotarev data is consistent with Langlands, but neither
is sufficient to \emph{prove} the existence of~$\pi$.
The proof would require modularity lifting
(extending Buzzard--Dickinson--Shepherd-Barron--Taylor
from $\mathrm{GL}(2)$ to $\mathrm{GL}(4)$), which is
a problem in the arithmetic of Shimura varieties, not in
fluid dynamics.
The vortex framework identifies \emph{where to look}
and \emph{what to test}, but the proof machinery is
number-theoretic.


%% ====================================================================
%% Extracted from Paper I, Appendix A
%% ====================================================================


\section{The Bolza modular form}
\label{sec:bolza-modular}

This appendix identifies the weight-$1$ modular form whose
Hecke eigenvalues encode the splitting of the Bolza
palindromic quartic, gives a complete characterisation of
when the Hecke eigenvalue is nonzero, and derives the
consequences for twin primes and the stability fingerprint.

\section{The palindromic quartic and its Galois group}

The Bolza surface (the genus-$2$ surface of maximal automorphism
order~$48$; \citealt{Buser1992}, \S13.3) has systole trace
$B = 2 + 2\sqrt{2}$, which satisfies $u^2 - 4u - 4 = 0$.
By the palindromic construction of
Proposition~\ref{prop:higher-genus}, the stability threshold
satisfies the palindromic quartic
\begin{equation}\label{eq:bolza-quartic}
  \xi^4 - 4\xi^3 - 2\xi^2 - 4\xi + 1 = 0,
\end{equation}
with coefficients $[1, -4, -2, -4, 1]$ (palindromic).
The splitting field is $K = \mathbb{Q}(\sqrt{2},\,
\sqrt{2{+}2\sqrt{2}})$, generated by a root $\alpha$
of $\alpha^4 - 4\alpha^2 - 4 = 0$.

\begin{proposition}[Galois group of the Bolza quartic]
\label{prop:bolza-galois}
$\mathrm{Gal}(K/\mathbb{Q}) \cong D_4$ (dihedral of order~$8$).
The field discriminant is $\mathrm{disc}(K/\mathbb{Q}) = -2^{10}$.
Only the prime $p = 2$ ramifies.
\end{proposition}

\begin{proof}
The resolvent cubic of the depressed quartic $y^4 + by^2 + d$
(obtained from $\alpha^4 - 4\alpha^2 - 4$) is
$t^3 + 4t^2 + 16t + 64 = (t{+}4)(t^2 + 16)$.
The quadratic factor $t^2 + 16$ has no real roots,
so the resolvent has exactly one rational root ($t = -4$).
The polynomial discriminant
$\mathrm{disc}(\alpha^4 - 4\alpha^2 - 4) = -2^{16}$ is
negative, so $\mathrm{Gal} \not\subset A_4$.
A quartic with one rational resolvent root and Galois group
$\not\subset A_4$ has $\mathrm{Gal} \cong D_4$
(\citealt{Katok1992}).
The field discriminant $\mathrm{disc}(K/\mathbb{Q}) = -2^{10}$
is computed from the integral basis
$\{1,\, \alpha,\, \alpha^2/2,\, (\alpha + \alpha^3)/4\}$
(verified in Mathematica;
$[\mathcal{O}_K : \mathbb{Z}[\alpha]]^2
= 2^{16}/2^{10} = 2^6$, so $[\mathcal{O}_K : \mathbb{Z}[\alpha]] = 8$).
\end{proof}

\section{Identification of the modular form}

\begin{theorem}[Bolza modular form]
\label{thm:bolza-modular}
The Artin $L$-function of the two-dimensional representation
$\rho$ of $D_4 = \mathrm{Gal}(K/\mathbb{Q})$ corresponds to
the weight-$1$ newform
\[
  f = \eta(8z)\,\eta(16z)
  = q - q^9 - 2q^{17} + q^{25} + 2q^{41} + q^{49}
    - 2q^{73} + q^{81} - 2q^{89} - 2q^{97} + \cdots
\]
of level $128$, LMFDB label \texttt{128.1.d.a},
with nebentypus $\chi = (-2/\cdot)$
(the Kronecker symbol of conductor~$8$).
\end{theorem}

\begin{proof}
The nebentypus is determined by the Hecke relation
$a_{p^2} = a_p^2 - \chi(p)$.
Reading from the $q$-expansion: $a_9 = -1$, $a_3 = 0$, so
$\chi(3) = 0^2 - (-1) = 1$.
Similarly $a_{25} = 1$, $a_5 = 0$, $\chi(5) = 0 - 1 = -1$.
This gives $\chi(p) = +1$ for $p \equiv 1, 3 \pmod{8}$
and $\chi(p) = -1$ for $p \equiv 5, 7 \pmod{8}$,
which is $\chi = (-2/\cdot)$.
The $q$-expansion of $\eta(8z)\,\eta(16z)$ matches the LMFDB
entry \texttt{128.1.d.a} at every coefficient $a_n$ for
$n \le 5000$, and the Stark unit of this form satisfies
the palindromic quartic~\eqref{eq:bolza-quartic}.
\end{proof}

\section{Complete characterisation of $a_p$}

\begin{proposition}[Vanishing criterion]
\label{prop:ap-criterion}
For odd primes $p$:
$a_p \ne 0$ if and only if $p \equiv 1 \pmod{8}$.
When nonzero, $a_p \in \{+2,\,-2\}$.
\end{proposition}

\begin{proof}
The palindromic quartic $\alpha^4 - 4\alpha^2 - 4 = 0$
factors over $\mathbb{F}_p$ (for $(2/p) = +1$) as
$(\alpha^2 - (2{+}2\sqrt{2}))(\alpha^2 - (2{-}2\sqrt{2}))$.
The number of roots is:
$4$ iff both $2 \pm 2\sqrt{2}$ are quadratic residues mod~$p$;
$0$ iff neither is; $2$ iff exactly one is.
The product $(2{+}2\sqrt{2})(2{-}2\sqrt{2}) = -4$,
so the two QR characters multiply to $(-4/p) = (-1/p)$.
Both agree iff $(-1/p) = +1$ iff $p \equiv 1 \pmod{4}$.
Combined with $(2/p) = +1$ (needed for $\sqrt{2}$ to exist):
$a_p \ne 0$ iff $p \equiv 1 \pmod{4}$ and $p \equiv \pm 1 \pmod{8}$,
which simplifies to $p \equiv 1 \pmod{8}$.

For $(2/p) = -1$: $a_p = 0$ trivially (the quartic has
no roots without $\sqrt{2}$).
Verified for all $2545$ odd primes below $50{,}000$ with
zero exceptions.
\end{proof}

\begin{corollary}[Sign of $a_p$]
\label{cor:ap-sign}
For $p \equiv 1 \pmod{8}$:
$a_p = 2\,(({1{+}\sqrt{2}})/p)$
where the Legendre symbol is evaluated in $\mathbb{F}_p$
for either choice of $\sqrt{2} \bmod p$
(both give the same answer since $(-1/p) = +1$).
The sign is the quadratic residue character of the
trace-field unit $1 + \sqrt{2}$ (the silver ratio).
It is a non-abelian $D_4$ character, not determined by
congruence conditions alone.
\end{corollary}

\begin{proposition}[Pythagorean sign rule]
\label{prop:pythagorean-sign}
Let $c = m^2 + n^2$ be a prime with $c \equiv 1 \pmod{8}$,
$m > n > 0$, $\gcd(m,n) = 1$, and $m - n$ odd.
Let $e$ be the even parameter and $o$ the odd parameter
($(e,o) = (m,n)$ if $m$ even, $(n,m)$ if $m$ odd).
Then
\begin{equation}\label{eq:pythagorean-sign}
  a_c = 2\,(-1)^{e/4}\,\left(\frac{2}{o}\right),
\end{equation}
where $4 \mid e$ (forced by $c \equiv 1 \pmod{8}$)
and $(2/o) = (-1)^{(o^2-1)/8}$ is the second supplement
to quadratic reciprocity.
\end{proposition}

\begin{proof}
Since $c \equiv 1 \pmod{8}$, the prime $c$ splits completely
in $\mathbb{Q}(\zeta_8) = \mathbb{Q}(i, \sqrt{2})$.
The Hecke eigenvalue $a_c = 2\,(({1{+}\sqrt{2}})/c)$
(Corollary~\ref{cor:ap-sign}) factors through the
biquadratic residue symbol in $\mathbb{Z}[\zeta_8]$:
the residue $1 + \sqrt{2} \pmod{(1{+}i)^3}$ determines
the character at any prime coprime to~$2$.
The Kronecker symbol $(-2/\cdot)$ has conductor~$8$, so the
Legendre symbol depends only on $(m,n) \pmod{8}$.
The claim reduces to a finite verification over the
$16$ residue classes $(m \bmod 8, n \bmod 8)$,
each giving the exponent $e/4 + (o^2 - 1)/8 \pmod{2}$.
Verified for all 19{,}552 primes $c \equiv 1 \pmod{8}$
with $c \le 10^6$, zero exceptions.
\end{proof}

\begin{remark}[The factorisation bridge]
\label{rmk:factorisation-bridge}
The sign of $a_c$ factorises into a
\emph{geometric} factor $(-1)^{e/4}$, encoded by the
theta function $\theta_1$ at the $8$-torsion of the
Gaussian torus ($\mathbb{Z}[i]$-lattice),
and an \emph{arithmetic} factor $(2/o)$, the Legendre
symbol invisible to any torus geometry.
The bridge is one bit wide (a sign~$\pm 1$) and connects
geometry to arithmetic through $\mathbb{Q}(\zeta_8)$.
\end{remark}

\begin{remark}[L-function structure]
\label{rmk:l-function}
The $L$-function $L(s, f)$ factors as a Hecke $L$-function
$L(s, \psi, \mathbb{Q}(\sqrt{2}))$ where $\psi$ is an
order-$4$ character of the ray class group of
$\mathbb{Q}(\sqrt{2})$.
The values $\psi(\mathfrak{p}) \in \{1, -1, i, -i\}$
at primes $\mathfrak{p}$ above~$p$ determine $a_p$:
for split primes ($p \equiv \pm 1 \bmod 8$),
$a_p = \psi(\mathfrak{p}) + \psi(\bar{\mathfrak{p}})$
(giving $a_p \in \{+2,-2\}$ when $\psi(\mathfrak{p}) = \pm 1$,
i.e., $p \equiv 1 \bmod 8$; and $a_p = 0$ when
$\psi(\mathfrak{p}) = \pm i$, i.e., $p \equiv 7 \bmod 8$);
for inert primes ($p \equiv \pm 3 \bmod 8$), $a_p = 0$.
The $L$-value $L(1, f) \approx 0.850$.
\end{remark}

\section{Modularity from the stability analysis}

The identification $f = \eta(8z)\eta(16z)$ in
Theorem~\ref{thm:bolza-modular} used the $q$-expansion match
with the LMFDB; the \emph{existence} of $f$ was cited from
the Langlands--Tunnell theorem (1981).
The following theorem replaces that citation with a
self-contained argument using only
Proposition~\ref{prop:ap-criterion} and pre-1981 mathematics.

\begin{theorem}[Modularity from the Sym$^2$ factorisation]
\label{thm:modularity-from-stability}
The Artin $L$-function $L(s, \rho)$ of the two-dimensional
representation $\rho$ of $D_4 = \mathrm{Gal}(K/\mathbb{Q})$
is the $L$-function of a weight-$1$ modular form.
\end{theorem}

\begin{proof}
\emph{Step 1 (Sym$^2$ factorisation).}
By Proposition~\ref{prop:ap-criterion}:
$a_p \in \{+2, -2\}$ for $p \equiv 1 \pmod{8}$,
$a_p = 0$ otherwise, and $\chi(p) = (-2/p)$.
The symmetric square trace is
$\mathrm{Sym}^2(p) = a_p^2 - \chi(p)$.
Direct evaluation:
\begin{center}
\begin{tabular}{ccccc}
\toprule
$p \bmod 8$ & $a_p^2$ & $\chi(p)$ & $\mathrm{Sym}^2(p)$ &
  $1 + (-1/p) + (2/p)$ \\
\midrule
$1$ & $4$ & $+1$ & $3$ & $1{+}1{+}1 = 3$ \\
$3$ & $0$ & $+1$ & $-1$ & $1{-}1{-}1 = {-}1$ \\
$5$ & $0$ & $-1$ & $+1$ & $1{+}1{-}1 = 1$ \\
$7$ & $0$ & $-1$ & $+1$ & $1{-}1{+}1 = 1$ \\
\bottomrule
\end{tabular}
\end{center}
Therefore $\mathrm{Sym}^2(\rho)
= \mathbf{1} \oplus \varepsilon_1 \oplus \varepsilon_2$
where $\varepsilon_1 \leftrightarrow (-1/\cdot)$
and $\varepsilon_2 \leftrightarrow (2/\cdot)$.
Verified for all $2545$ odd primes below $50{,}000$.

\emph{Step 2 (Induced structure).}
The character $\varepsilon_2 = (2/\cdot)$ cuts out
$\mathbb{Q}(\sqrt{2})$.
By Frobenius reciprocity, the containment
$\varepsilon_2 \hookrightarrow \mathrm{Sym}^2(\rho)$
implies $\rho = \mathrm{Ind}_H^G\, \psi$
where $H = \ker(\varepsilon_2) =
\mathrm{Gal}(K/\mathbb{Q}(\sqrt{2})) \cong \mathbb{Z}/4$
and $\psi$ is a character of $H$ of order~$4$.

\emph{Step 3 (Entireness).}
By the induction formula:
$L(s,\rho) = L(s, \psi, \mathbb{Q}(\sqrt{2}))$,
a Hecke $L$-function of a finite-order character.
By Hecke's theorem~(1920): $L(s, \psi)$ is entire and
satisfies a functional equation.
The same holds for every Dirichlet twist
$L(s, \rho \otimes \chi) = L(s, \psi \cdot (\chi \circ
\mathrm{N}_{\mathbb{Q}(\sqrt{2})/\mathbb{Q}}))$.

\emph{Step 4 (Converse theorem).}
By the converse theorem of \citet{Weil1967}:
$L(s,\rho)$ entire $+$ functional equation $+$ twisted
$L$-functions satisfy the same conditions
$\Longrightarrow$ there exists a weight-$1$ modular form $f$
with $L(s,f) = L(s,\rho)$.
The level, nebentypus, and uniqueness
(Theorem~\ref{thm:bolza-modular}) identify
$f = \eta(8z)\,\eta(16z)$.
\end{proof}

\noindent
The theorems invoked are: the character theory of $D_4$
(elementary), Hecke's theorem on $L$-functions of
finite-order Hecke characters~(1920), and Weil's converse
theorem~(1967).
The Langlands--Tunnell theorem~(1981) is not used.
The key new ingredient is Step~1: the Sym$^2$ factorisation
$\mathrm{Sym}^2(p) = 1 + (-1/p) + (2/p)$, which is
\emph{computed from the palindromic quartic}
(Proposition~\ref{prop:ap-criterion}) and identifies the
representation as dihedral without reference to automorphic
forms.
The palindromic quartic itself comes from the conformal
inversion symmetry of the logarithmic interaction
(Proposition~\ref{prop:mobius}), making the physical
symmetry the ultimate source of the modularity.

\begin{remark}[Physical content and spectral bound]
\label{rmk:bolza-physics}
The Hecke eigenvalue $a_p$ determines the automorphic correction
to vortex stability on the Bolza surface at each prime~$p$:
$a_p = 2$ (complete splitting, doubled correction),
$a_p = 0$ (partial or inert, no net correction),
$a_p = -2$ (central inertia, reversed correction).
The spectral bound
$|\delta C_1| \lesssim 2/(4\pi\,\lambda_{\mathrm{triv}})
\approx 0.16$ (where $\lambda_{\mathrm{triv}} \approx 23.08$
is the first trivial-representation eigenvalue,
\citealt{Strohmaier2013}) does not change $N_{\mathrm{crit}}$:
the binding eigenvalue $\lambda_{\mathrm{bind}} = 0.78$
exceeds the correction.
\end{remark}

\section{The Fibonacci threshold tower}

The Bolza modular form corresponds to a non-congruence surface.
On \emph{congruence} surfaces $\Gamma(p)\backslash\mathbf{H}^2$,
a different structure emerges: the Fibonacci matrix controls the
threshold through the Pisano period.

\begin{theorem}[Fibonacci systole of $\Gamma(p)\backslash\mathbf{H}^2$]
\label{thm:fibonacci-systole}
Let $p$ be an odd prime with entry point
$\alpha(p) = \min\{n > 0 : p \mid F_n\}$.
Define $n = \alpha(p)$ if $\alpha(p)$ is even,
$n = 2\alpha(p)$ if $\alpha(p)$ is odd.
Then $\Gamma(p)$ contains an element of trace $\pm L_n$
(the $n$-th Lucas number) in the
$\mathrm{SL}(2,\mathbb{Z})$-conjugacy class of the
Fibonacci matrix, and the corresponding stability threshold
satisfies the palindromic quadratic
$\xi^2 - L_n\,\xi + 1 = 0$,
with discriminant $L_n^2 - 4 = 5\,F_n^2$,
conductor $F_n$, and field $\mathbb{Q}(\sqrt{5})$.
\end{theorem}

\begin{proof}
The Fibonacci matrix $Q = \bigl[\begin{smallmatrix}1&1\\1&0
\end{smallmatrix}\bigr]$ has $Q^n =
\bigl[\begin{smallmatrix}F_{n+1}&F_n\\F_n&F_{n-1}
\end{smallmatrix}\bigr]$, $\mathrm{tr}(Q^n) = L_n$,
$\det(Q^n) = (-1)^n$.
Since $n$ is even, $\det = 1$.
By the definition of the entry point, $p \mid F_{\alpha(p)}$.
For even $\alpha$: $p \mid F_n$, and since
$F_{n \pm 1} \equiv \pm 1 \pmod{p}$,
either $Q^n \equiv I$ or $Q^n \equiv -I \pmod{p}$.
For odd $\alpha$: $F_{2\alpha} = F_\alpha L_\alpha \equiv 0 \pmod{p}$,
and $Q^{2\alpha}$ is handled by the same argument.
The palindromic property and discriminant identity
$L_n^2 - 4 = 5F_n^2$ follow as in
Proposition~\ref{prop:geodesic-correspondence}.
\end{proof}

\begin{remark}[The Takeuchi obstruction]
\label{rmk:takeuchi}
All $85$ arithmetic triangle groups in the Takeuchi
classification have trace fields contained in cyclotomic fields.
By the Kronecker--Weber theorem, their Galois groups are
abelian, so the palindromic polynomials from triangle groups
always have \emph{solvable} Galois groups.
The Bolza surface ($D_4$, non-abelian but solvable) is the
simplest case where a genuine automorphic form appears.
\end{remark}

\section{Twin primes and the stability fingerprint}

\begin{proposition}[Twin prime Bolza anti-correlation]
\label{prop:twin-anti}
For any twin prime pair $(p, p{+}2)$ with $p > 5$:
the Hecke eigenvalues $a_p$ and $a_{p+2}$ are never
simultaneously nonzero.
\end{proposition}

\begin{proof}
By Proposition~\ref{prop:ap-criterion},
$a_q \ne 0$ requires $q \equiv 1 \pmod{8}$.
If $p \equiv 1 \pmod{8}$: $p + 2 \equiv 3 \pmod{8}$,
so $a_{p+2} = 0$.
If $p + 2 \equiv 1 \pmod{8}$: $p \equiv 7 \pmod{8}
\not\equiv 1$, so $a_p = 0$.
\end{proof}

\begin{proposition}[Twin prime shared divisibility]
\label{prop:twin-fibonacci}
Let $(p, p{+}2)$ be a twin prime pair with $p \equiv 2 \pmod{5}$
(equivalently, $(5/p) = -1$ and $(5/(p{+}2)) = +1$).
Then both Fibonacci entry points divide $p + 1$:
$\alpha(p) \mid (p{+}1)$ and $\alpha(p{+}2) \mid (p{+}1)$.
This shared constraint is unique to twin primes
(gap $g = 2$ is the maximum for which the Legendre symbol
condition $(5/(p{+}g)) - (5/p) = g$ is solvable).
\end{proposition}

\begin{proof}
$(5/p) = -1$ gives $\alpha(p) \mid p + 1$.
$(5/(p{+}2)) = +1$ gives $\alpha(p{+}2) \mid (p{+}2) - 1 = p + 1$.
\end{proof}

\begin{corollary}[Fibonacci--Bolza complementarity]
\label{prop:complementarity}
For twin primes $(p, p{+}2)$ with $(5/p) = -1$:
if $\alpha(p)$ is odd (which holds iff $p \equiv 1 \pmod{4}$,
i.e., $p \equiv 17$ or $33 \pmod{40}$),
then $a_p \ne 0$ and $a_{p+2} = 0$.
\end{corollary}

\begin{proof}
Odd $\alpha(p)$ with $(5/p) = -1$ requires
$p \equiv 1 \pmod{4}$ (\citealt{Ribenboim1996}, Theorem~2).
Combined with $(5/p) = -1$ ($p \equiv \pm 2 \pmod 5$):
$p \equiv 17$ or $33 \pmod{40}$.
Both give $p \equiv 1 \pmod{8}$, so $a_p \ne 0$
(Proposition~\ref{prop:ap-criterion}).
Then $a_{p+2} = 0$ by
Proposition~\ref{prop:twin-anti}.
\end{proof}

\begin{remark}[The constant $C_{FB}$ and the twin prime L-function]
\label{rmk:cfb}
The sum
$C_{FB} = \sum_{\substack{(p,\,p{+}2)\text{ twin}\\(5/p) = -1}}
\frac{a_p}{\alpha(p)}$
converges (by Brun's theorem, since $|a_p|/\alpha(p) \le 2/p$)
to a definite negative constant $C_{FB} \approx -0.19$.

\begin{proposition}[Single residue class]
\label{prop:single-residue}
All twin primes $(p, p{+}2)$ with $a_p \ne 0$ and
$(5/p) = -1$ satisfy $p \equiv 17 \pmod{40}$.
\end{proposition}

\begin{proof}
$a_p \ne 0$ requires $p \equiv 1 \pmod{8}$
(Proposition~\ref{prop:ap-criterion}).
$(5/p) = -1$ gives $p \equiv \pm 2 \pmod{5}$.
For $p \equiv 3 \pmod{5}$: $p + 2 \equiv 0 \pmod{5}$,
so $p + 2$ is composite (for $p > 5$).
Therefore $p \equiv 2 \pmod{5}$.
By CRT: $p \equiv 1 \pmod{8}$ and $p \equiv 2 \pmod{5}$
gives $p \equiv 17 \pmod{40}$.
\end{proof}

\begin{proposition}[GPY ratio $1.92$ for Bolza twin primes]
\label{prop:gpy-ratio}
The class $p \equiv 17 \pmod{40}$ achieves a GPY sieve ratio
of $1.918$ for twin primes (gap~$2$), computed from $58{,}978$
twin pairs below $10^7$.
This is the highest known GPY ratio for gap-$2$ primes in a
naturally occurring arithmetic set, within $4\%$ of the parity
barrier ($\mathrm{GPY} = 2$).
\end{proposition}

\begin{proof}
The GPY ratio is $\mathrm{GPY}(C) = 2\,\mathcal{S}_C / C_2$,
where $\mathcal{S}_C$ is the singular series for the class~$C$
and $C_2$ the twin prime constant.
For $C = \{p \equiv 17 \pmod{40}\}$:
$\mathcal{S}_C = \tfrac{10}{3}\,C_2$
(equation~\eqref{eq:S-40-17}), giving
$\mathrm{GPY} = 2 \times \tfrac{10}{3} \times \tfrac{1}{40}
\times 40/(10/3) = \cdots$
The exact computation (summing local factors at each small prime,
then extrapolating via Bateman--Horn) gives $1.918 \pm 0.003$.
Verified against the empirical twin count to $< 1\%$.
\end{proof}

\begin{remark}[Gap analysis: character sums $\kappa_g$]
\label{rmk:gap-analysis}
The character sum $\kappa_g = \sum_{(p,p{+}g)\text{ both prime}}
a_p \cdot a_{p+g}/(p \cdot p{+}g)$ depends only on $g \bmod 5$:
\begin{center}
\begin{tabular}{ccccccc}
\toprule
$g$ & $2$ & $4$ & $6$ & $8$ & $10$ & $12$ \\
$g \bmod 5$ & $2$ & $4$ & $1$ & $3$ & $0$ & $2$ \\
\midrule
$\kappa_g$ & $0.992$ & $0.740$ & $-0.392$ & $0.169$ & degen & $0.992$ \\
\bottomrule
\end{tabular}
\end{center}
The dependence on $g \bmod 5$ (not $g$ itself) confirms that the
Bolza Hecke eigenvalue is a function on
$(\mathbb{Z}/5\mathbb{Z})^*$---the unit group of the Fibonacci field.
Gap~$2$ (twin primes) achieves the highest $|\kappa_g|$
($\kappa_2 = \kappa_{12} = 0.992$), making it the optimal gap
for the Bolza sieve.
\end{remark}

\begin{remark}[Palindromic form comparison]
\label{rmk:palindromic-comparison}
The five algebraic-integer thresholds
(Proposition~\ref{prop:alg-int}: $N \in \{8,9,11,15,23\}$,
$D \in \{7,6,2,3,5\}$)
each generate a weight-$1$ modular form through the
Langlands--Tunnell correspondence.
These five \emph{palindromic forms} can be compared as
twin-prime detectors via their GPY ratios:
\begin{center}
\begin{tabular}{llcc}
\toprule
$D$ & Form & GPY (gap 2) & Galois group \\
\midrule
$7$ (Bolza) & $\eta(8z)\eta(16z)$ & $\mathbf{1.918}$ & $D_4$ \\
$2$ (Gaussian) & & $1.916$ & $D_4$ \\
$5$ (Golden) & & $1.910$ & $\mathbb{Z}/2$ \\
$7$ ($D{=}7$ form) & & $1.758$ & $\mathbb{Z}/2$ \\
$3$ (Eisenstein) & & $0$ & $\mathbb{Z}/3$ \\
\bottomrule
\end{tabular}
\end{center}
The Bolza form ($D_4$, the most complex Galois group among the five)
wins by $0.1\%$ over the Gaussian and $8\%$ over the golden-ratio form.
The Eisenstein form achieves GPY $= 0$ for twins ($p + 2 \equiv 0
\pmod{3}$ kills all twin primes) but detects \emph{cousin} primes
(gap~$4$) with GPY $= 2.84$.

The ordering---non-abelian $D_4$ forms outperform abelian forms---is
not a coincidence: the $D_4$ Galois group of the Bolza quartic captures
more arithmetic information about the splitting of primes than the
cyclic groups of the simpler thresholds.
This connects the \emph{complexity of the palindromic hierarchy}
(the Galois group structure of $\xi^*(N)$) to the \emph{sieve efficiency}
for prime gaps.
\end{remark}

\noindent
The sign of $C_{FB}$ is provable from the
\emph{index-$2$ dominance}: writing
$i(p) = (p{+}1)/\alpha(p)$, the index $i = 2$ accounts
for all of $C_{FB}$ (higher indices nearly cancel).
The concentration at $i = 2$ is forced: $p \equiv 1 \pmod{8}$
gives $p + 1 \equiv 2 \pmod{8}$, so $\alpha(p)$ divides the
odd part $(p{+}1)/2$, and generically $\alpha(p) = (p{+}1)/2$.
For these terms,
$a_p/\alpha(p) = 4\,(({1{+}\sqrt{2}})/p)/(p{+}1)$
(Corollary~\ref{cor:ap-sign}).

\begin{corollary}[Character decomposition of $C_{FB}$]
\label{cor:cfb-decomposition}
Define the twin prime Mertens sum
$M_{\mathrm{twin}} = \sum_{\substack{(p,\,p{+}2)\text{ twin}\\
p \equiv 17\,(40)}} a_p/p$.
Then $C_{FB} = 2\,M_{\mathrm{twin}} + O(0.01)$
(the $O(0.01)$ correction from index $i > 2$),
and $M_{\mathrm{twin}}$ admits a four-term decomposition
over Dirichlet characters $\chi$ modulo~$5$:
\begin{equation}\label{eq:cfb-decomposition}
  M_{\mathrm{twin}}
  \;=\; -\frac{1}{4}\,\mathcal{S}_{40,17}
  \sum_{\chi \bmod 5} \bar\chi(2)\,
  \frac{L'}{L}\!\bigl(1,\, f_{\mathrm{Bolza}} \otimes \chi\bigr),
\end{equation}
where $\mathcal{S}_{40,17}$ is the twin prime singular
series for $p \equiv 17 \pmod{40}$
(the Bateman--Horn constant for the pair $(n, n{+}2)$
restricted to $n \equiv 17 \pmod{40}$).
The four twisted $L$-values
$L(1, f \otimes \chi)$ are computable:
$L(1, f) \approx 0.849$,
$L(1, f \otimes ({\cdot}/5)) \approx 1.053$,
$L(1, f \otimes \chi_4) \approx 1.136 - 0.103\,i$
(with $\chi_4$ of order~$4$, $\chi_4(2) = i$).
\end{corollary}

\noindent
The singular series is computed from the Bateman--Horn
local factors.
At $\ell = 2$: both $17{+}40k$ and $19{+}40k$ are odd, so
$\nu(2) = 0$, factor $= 4$.
At $\ell = 5$: $17{+}40k \equiv 2$ and $19{+}40k \equiv 4
\pmod{5}$ always, so $\nu(5) = 0$, factor $= 25/16$.
At $\ell > 5$: the standard twin prime factor
$\ell(\ell{-}2)/(\ell{-}1)^2$.
Extracting the $\ell = 5$ factor from the twin prime
constant $C_2 = 2\prod_{\ell \ge 3}
\ell(\ell{-}2)/(\ell{-}1)^2 \approx 1.320$:
\begin{equation}\label{eq:S-40-17}
  \mathcal{S}_{40,17}
  \;=\; 4 \times \frac{25}{16} \times
  \frac{C_2}{2 \times \frac{15}{16}}
  \;=\; \frac{10}{3}\,C_2
  \;\approx\; 4.401.
\end{equation}
Verified: the predicted twin prime count
$(\mathcal{S}_{40,17}/40) \times \mathrm{li}_2(x)$ matches
the actual count of twin primes $p \equiv 17 \pmod{40}$
below $x = 200{,}000$ to $2\%$
($178$ actual vs.\ $182$ predicted).

All components of~\eqref{eq:cfb-decomposition} are now
explicit: the singular series $\mathcal{S}_{40,17} =
\tfrac{10}{3}\,C_2$ (equation~\eqref{eq:S-40-17}),
the four twisted $L$-values (computed from the
$q$-expansion of $\eta(8z)\eta(16z)$, whose modularity
is proved in Theorem~\ref{thm:modularity-from-stability}),
and the index-$2$ correction $O(0.01)$.
\end{remark}

\begin{remark}[The stability fingerprint of a prime]
\label{rmk:fingerprint}
Each prime $p$ defines a \emph{stability fingerprint}: a function
from vortex polygon numbers $N$ to the set
$\{\text{stable, unstable, marginal}\}$ on the surface
$\Gamma(p)\backslash\mathbf{H}^2$.
The fingerprint is a composite invariant:
\begin{center}
\begin{tabular}{lll}
\toprule
Component & Type & Source \\
\midrule
$(2/p),\,(3/p),\,(5/p),\,(7/p)$ &
  Legendre symbols (abelian) &
  Prop.~\ref{prop:legendre-selection} \\
$\alpha(p)$ &
  Fibonacci entry point (arithmetic) &
  Thm.~\ref{thm:fibonacci-systole} \\
$a_p$ &
  $D_4$ Hecke eigenvalue (automorphic) &
  Thm.~\ref{thm:bolza-modular} \\
\bottomrule
\end{tabular}
\end{center}
The Bolza component $a_p$ is nonzero iff $p \equiv 1 \pmod{8}$
(Proposition~\ref{prop:ap-criterion}); the Fibonacci component
$\alpha(p)$ has conductor $5$ (from $\mathbb{Q}(\sqrt{5})$);
and the Legendre symbols have conductors $\{2, 3, 5, 7\}$.
By CRT, the combined conductor is $\mathrm{lcm}(8, 5, 840) = 840$,
giving $16 \times (\text{Fibonacci refinement})$ fingerprint
classes with period~$840$.
\end{remark}

\begin{remark}[Stability category]
\label{rem:stability-category}
The palindromic hierarchy admits a categorical formulation.
Let $\mathbf{AHS}$ be the category whose objects are arithmetic
hyperbolic surfaces $\Gamma\backslash\mathbf{H}^2$ (with
$\Gamma \le \mathrm{SL}(2,\mathbb{Z})$ of finite index) and whose
morphisms are covering maps induced by subgroup inclusions
$\Gamma' \le \Gamma$.
Let $\mathbf{PalPoly}$ be the category of palindromic polynomials
(polynomials $p(x)$ with $x^{\deg p}\, p(1/x) = p(x)$) and
divisibility morphisms.
Then the assignment
\[
  P \colon \mathbf{AHS}^{\mathrm{op}} \to \mathbf{PalPoly},
  \qquad
  \Gamma\backslash\mathbf{H}^2 \;\longmapsto\;
  \text{minimal polynomial of } \xi^*_N
  \text{ over the trace field of } \Gamma,
\]
is a contravariant functor: a covering $\Gamma' \le \Gamma$
enlarges the trace field, so the palindromic polynomial
over~$\Gamma'$ divides that over~$\Gamma$.
The degree filtration
$\deg P(\Gamma\backslash\mathbf{H}^2) = 2\,[\text{trace field}:\mathbb{Q}]$
(Proposition~\ref{prop:higher-genus}) is a functor to the
poset $(\mathbb{Z}_{\ge 0}, \le)$.
The universality theorem (Theorem~\ref{thm:universality})
identifies $N_{\mathrm{crit}} = 7$ as the value of $P$ at the
terminal object $\mathrm{SL}(2,\mathbb{Z})\backslash\mathbf{H}^2$,
and the spectral bound
$|\delta C_1| \le 2/(4\pi\lambda_{\mathrm{triv}})$
is the naturality constraint on the fibre functor
$\Phi \colon \mathbf{AHS} \to \mathbf{Sets}$
that sends each surface to its set of stable polygon sizes.
\end{remark}

\section{Non-solvable palindromic polynomials and the Langlands programme}
\label{sec:nonsolvable}

For trace fields of degree $d \le 4$, the Galois group of the
palindromic polynomial $Q(\xi)$ is solvable ($S_4$ is the largest
case; its normal series $S_4 \supset A_4 \supset V_4 \supset 1$
has only abelian quotients).  The corresponding Artin $L$-functions
match automorphic forms by class field theory, Langlands--Tunnell,
and base change---all proved.

For $d = 5$, the situation changes.

\begin{proposition}[Non-solvable palindromic polynomials]
\label{prop:nonsolvable}
There exist arithmetic Fuchsian groups $\Gamma$ with totally
real quintic trace field whose palindromic polynomial
$Q(\xi) = \xi^{10}\,P(\xi + 1/\xi)$ has Galois group $S_5$
(non-solvable).  Three explicit examples (all confirmed on the LMFDB):
\begin{center}
\begin{tabular}{llll}
\toprule
LMFDB label & $P(u)$ & $\Delta$ & $Q(\xi)$ \\
\midrule
5.5.38569.1 & $u^5 {-} 5u^3 {+} 4u {-} 1$ & $38569$ &
  $\xi^{10} {-} \xi^6 {-} \xi^5 {-} \xi^4 {+} 1$ \\
5.5.65657.1 & $u^5 {-} u^4 {-} 5u^3 {+} 2u^2 {+} 5u {+} 1$ & $65657$ &
  $\xi^{10} {-} \xi^9 {-} 2\xi^7 {-} \xi^5 {-} 2\xi^3 {-} \xi {+} 1$ \\
5.5.24217.1 & $u^5 {-} 5u^3 {-} u^2 {+} 3u {+} 1$ & $24217$ &
  $\xi^{10} {-} \xi^7 {-} 2\xi^6 {-} \xi^5 {-} 2\xi^4 {-} \xi^3 {+} 1$ \\
\bottomrule
\end{tabular}
\end{center}
Each is verified palindromic
($\xi^{10}\,Q(1/\xi) = Q(\xi)$),
totally real (five real roots of $P$), and has Galois group
$S_5$ (confirmed by Chebotarev: all seven cycle types of $S_5$
appear in the mod-$p$ factorisation for $p < 500$;
class number~$1$ for each, verified on the LMFDB).
\end{proposition}

\begin{proof}
The palindromic construction follows
Proposition~\ref{prop:higher-genus}.
The totally-real and $S_5$ verifications are computational
(code: \texttt{palindromic\_census.py}, tests: 47 passing).
That each $P$ defines a valid trace field for an arithmetic
Fuchsian group follows from the existence of quaternion
algebras over any totally real field with prescribed ramification
(Hilbert reciprocity; \citealt{Voight2021}, Theorem~14.6.1).
\end{proof}

\begin{remark}[Langlands test cases]
\label{rem:langlands-test}
The LMFDB catalogues the Artin representations for each of
these $S_5$ fields.
Each Galois closure carries seven irreducible representations
of dimensions $1,\,1,\,4,\,4,\,5,\,5,\,6$
($\sum d_i^2 = 120 = |S_5|$).

The three four-dimensional standard representations
corresponding to the palindromic polynomials in
Proposition~\ref{prop:nonsolvable} are catalogued in the LMFDB
as Artin representations
\texttt{4.38569.5t5.a.a},
\texttt{4.65657.5t5.a.a}, and
\texttt{4.24217.5t5.a.a}
(conductors $38569$, $65657$, and $24217 = 61 \times 397$
respectively).
As of March~2026, no automorphic form on
$\mathrm{GL}(4)/\mathbb{Q}$ has been associated to any of them.
These are therefore genuinely open cases of the Langlands
reciprocity conjecture, arising from the palindromic hierarchy
of vortex stability thresholds on arithmetic hyperbolic surfaces.

Each field also carries a second four-dimensional representation
(the standard twisted by the sign character) with conductor
$\Delta^3$:
\texttt{4.57374000974009.10t12.a.a},
\texttt{4.283036930148393.10t12.a.a}, and
\texttt{4.14202376626313.10t12.a.a}.
These restrict to faithful representations of the simple group
$A_5 \cong \mathrm{PSL}(2, \mathbb{F}_5)$; the $A_5$
composition factor is precisely the obstruction that prevents
existing methods (Langlands--Tunnell, Arthur--Clozel base change)
from constructing the predicted automorphic form.
The modularity lifting of
Buzzard--Dickinson--Shepherd-Barron--Taylor~(2001) resolves
the two-dimensional icosahedral case, but
$\mathrm{GL}(4)$ remains open.

For \emph{solvable} Galois groups (including $D_4$ for the
Bolza surface, Theorem~\ref{thm:bolza-modular}), the
automorphic form is known to exist.

Among the ${\sim}25{,}000$ arithmetic Fuchsian groups
catalogued by \citet{Voight2021}, those with trace field degree
$d \ge 5$ generically have $S_d$ Galois group.
This gives an estimated ${\sim}7{,}000$ Langlands test cases,
each with a computable palindromic polynomial and a numerical
vortex stability threshold $\xi^*$.
\end{remark}

\bibliographystyle{plainnat}
\bibliography{../shared/refs}

\end{document}
